\section{Size-specific computations and bounded error}\label{sec:robustness}
Restriction compares structures of different sizes. In this section we keep
the domain fixed and allow the target to be chosen separately for each
domain size. No new finite structure needs to be built, because of the way
in which a deletion acted in the first place: by
Section~\ref{sec:responses}, a deletion affected a composition only through
the three responses, that is, only by removing a color from the palette
that the hub can see. Recoloring the hub's edges into $\cell\lambda$ with
some other color has exactly the same effect and removes no vertex. The
same three responses therefore control such a recoloring of $\M$, on a
domain that never changes.

\subsection{Transport of type values on a fixed domain}
For a map $f:\Col\to\Col$, we change only the colors of the edges at the
hub:
\[
 \colr_f(h,u)=\colr_f(u,h)=f(\lambda)\quad(u\in\cell\lambda),\qquad
 \colr_f(u,u')=\colr(u,u')\quad(u\ne u'\in\cell\Col),
\]
and put $J=f(\Col)$. Write $\ctype\lambda^f$ for the new color types and
$\tp_f$ for the new typing map; the anchor types and $\tpeq$ are unchanged.
Every color still occurs on edges between cell vertices, so these types are
nonempty and partition $D^2$. What is lost is homogeneity: seen from the
hub, only the colors in $J$ now occur, whereas every color still occurs
elsewhere. Hence Lemma~\ref{lem:profile} fails for $\colr_f$, and the
type-constant arrays with respect to the recolored types need not be closed
under aggregation. The lemma below replaces this closure property.

For a type-constant $P$, let $f_\ast P$ be the array that takes the value
$P(\ctype\lambda)$ on $\ctype\lambda^f$ and agrees with $P$ on all other
types. This map commutes with uniform pointwise functions, transpose, and
the permitted constants. In particular, it does not change the values at
the read-off pairs:
\begin{equation}\label{eq:transport-anchor}
 (f_\ast P)(x,h)=P(x,h)\qquad(x\in X).
\end{equation}
The map $f_\ast$ transports values between types; it need not preserve
individual entries elsewhere.

The three responses now have to be compatible with $f_\ast$ rather than
with restriction.
\begin{lemma}[Three-response recoloring]\label{lem:recolor}
For type-constant $P,Q$, if the responses \eqref{eq:responses} agree on
$J=f(\Col)$ and on $\Col$, then
\[
 f_\ast(P\agg\odot Q)=(f_\ast P)\agg\odot(f_\ast Q).
\]
\end{lemma}
\begin{proof}
Both sides are arrays on the same domain $D$; we compare them at a pair
$(u,v)$. We call $\tau$ the \emph{name} of the recolored type of $(u,v)$: it
is $\ctype\lambda$ when $\tp_f(u,v)=\ctype\lambda^f$, and it is the type
itself for the four kinds of types that $\colr_f$ leaves unchanged, namely
$\tpeq$, the $\tpl x$, the $\tpr x$ and the $\tpa x{x'}$. By the definition of $f_\ast$ the
left-hand side at $(u,v)$ equals $(P\agg\odot Q)(\tau)$, which by
\eqref{eq:profile-value} is the join of $P(\tau_1)\odot Q(\tau_2)$ over the
reference profiles $(\tau_1,\tau_2)\in W_\tau$. The right-hand side at
$(u,v)$ is the join of the same expression over the names of the recolored
profiles,
\[
 W^f_{(u,v)}=\bigl\{\text{names of }
 \bigl(\tp_f(u,z),\tp_f(z,v)\bigr):z\in D\bigr\},
\]
because $(f_\ast P)(u,z)=P(\tau_1)$ when $\tau_1$ is the name of
$\tp_f(u,z)$, and likewise for $Q$. So it suffices to show that the two
joins agree. Since $\colr_f$ differs from $\colr$ only on the edges at the
hub, and there replaces the palette $\Col$ by $J$, we sort the witnesses
into three groups.

\emph{Anchor witnesses and endpoint witnesses.} An anchor witness $z$ gives
two types among the anchor types, which $\colr_f$ does not touch, so it
contributes its reference term. A witness $z\in\{u,v\}\cap U$ gives
$(\tpeq,\tau)$ or $(\tau,\tpeq)$ if $u\ne v$, and $(\tpeq,\tpeq)$ if $u=v$;
these are the reference profiles that $W_\tau$ attaches to the endpoints of
a pair of type $\tau$, because $\tau$ is by construction the name of the
recolored type of $(u,v)$.

\emph{The hub as a witness, when neither endpoint is the hub.} The hub then
contributes exactly one profile name, and that name already occurs among the
cell witnesses at the same pair, so idempotence of the join absorbs it. In
detail, writing $\cell\mu$ and $\cell\nu$ for the cells of the left and the
right endpoint when these are cell vertices: at two distinct cell endpoints
the hub gives
$(\ctype{f(\mu)},\ctype{f(\nu)})$, and Lemma~\ref{lem:palette}(a) realizes
every ordered pair of colors there; at equal cell endpoints it gives a
paired name $(\ctype{f(\mu)},\ctype{f(\mu)})$, and Lemma~\ref{lem:palette}(b)
realizes every paired color; at an anchor--cell pair it gives $(\tpl{x},
\ctype{f(\nu)})$ or $(\ctype{f(\mu)},\tpr{x'})$, and
Lemma~\ref{lem:palette}(b) realizes every color on the cell side; and at a
pair of two anchors it gives $(\tpl{x},\tpr{x'})$, which every core witness
gives.

\emph{Cell witnesses.} These are the entries of the table
\eqref{eq:response-table}, read with $I=J$ and with the cell endpoints
ranging over all of $\cell\Col$. Indeed a witness $z\in\cell\mu$ is now seen
\emph{from the hub} through the name $\ctype{f(\mu)}$, so as $\mu$ runs over
$\Col$ that name runs over $J$; at any other endpoint every color is still
available, by Lemma~\ref{lem:palette}. The entries that depend on $I$ are
$r_I(p_x)$, $\ell_I(q_{x'})$, $d_I$, $\bigvee_\mu\ell_I(q_\mu)$ and
$\bigvee_\mu r_I(p_\mu)$, and the hypothesis turns each of them into its
value at $I=\Col$, which is the reference entry. The remaining entries do
not depend on $I$.

The three groups exhaust $D$, and together they reproduce the join over
$W_\tau$ at every pair $(u,v)$. Beyond the idempotence of the join, no law
of $\odot$ is used.
\end{proof}

\needspace{9\baselineskip}
As in Section~\ref{sec:simultaneous}, choosing a support at every gate
handles an entire computation at once.
\begin{theorem}[Simultaneous transport]\label{thm:transport}
For a computation $\Comp$ with type-constant inputs on $D$, there is
$K\subseteq\Col$ with $|K|\le\wid{\Comp}$ such that every map
$f:\Col\to\Col$ with $f(\Col)\supseteq K$ transports every intermediate
value by $f_\ast$: the value of each gate on the inputs $f_\ast P$ is
$f_\ast$ of its value on the original inputs $P$. One may use the same
support as in Theorem~\ref{thm:preservation}.
\end{theorem}
\begin{proof}
Use the support selected at each aggregation gate on the reference
computation. Free
operations commute with $f_\ast$ and Lemma~\ref{lem:recolor} handles
aggregation, so induction in evaluation order gives the claim. If the
support number is infinite, use $K=\Col$ and the same argument, since
$f(\Col)=\Col$ preserves every response. Each shared gate is charged once.
\end{proof}
All comparisons now take place on one and the same domain $D$. The wiring,
kernels, scalar constants, and pointwise functions may therefore have been
chosen for that domain size, provided that they are fixed before the input
is given and are uniform across entries. Besides the inputs, the arrays
available as primitives remain those listed in
Section~\ref{sec:responses}, namely the constant arrays and the equality
array. What is ruled out is any further primitive array whose $(u,v)$ entry
varies with $u$ and $v$, such as an array recording the cell of each
vertex, since such an array would give the computation exactly the
information that a type hides.

\subsection{Detectable labels at a fixed size}
For $f:\Col\to\Col$, let $\M_f$ be the structure on $D$ whose relations are
$f_\ast R^{\M}$, that is, the reference structure with its hub edges
recolored. Suppose that a Boolean-output query $\varphi$ has a
type-constant interpretation on $\M$ together with a set $\Det\subseteq\Col$
of labels, with corresponding anchors $x_\lambda$, that are
\emph{detectable under recoloring} in the following sense: for every
$\lambda\in\Det$,
\begin{equation}\label{eq:transport-decoder}
 \varphi^{\M_f}(x_\lambda,h)\ \text{holds}\iff\lambda\in f(\Col).
\end{equation}
For a non-Boolean $S$, the two truth values are represented by two fixed
distinct output states.

For each $\lambda\in\Det$, fix a color $\lambda'\ne\lambda$ and let
$f_\lambda:\Col\to\Col$ be the map that fixes all other colors and sends
$\lambda$ to $\lambda'$; then $f_\lambda(\Col)=\Col\setminus\{\lambda\}$,
and $\M_{f_\lambda}$ is $\M$ with the hub edges into $\cell\lambda$
recolored $\lambda'$. All of these structures have the same domain $D$, and
they are the fixed-size counterparts of the deletions used in
Corollary~\ref{cor:detectable}.
\begin{proposition}[A fixed-size test]\label{prop:fixed-test}
Under these assumptions, any aggregation computation $\Comp$ agreeing with
$\varphi$ on $\M$ and on the $|\Det|$ structures $\M_{f_\lambda}$,
$\lambda\in\Det$, has $\wid{\Comp}\ge|\Det|$.
\end{proposition}
\begin{proof}
If $\wid{\Comp}<|\Det|$, a support from Theorem~\ref{thm:transport} omits
some $\lambda\in\Det$. The structures $\M$ and $\M_{f_\lambda}$ then give
the same output at $(x_\lambda,h)$ by \eqref{eq:transport-anchor}, but
opposite correct answers by \eqref{eq:transport-decoder}. These $|\Det|+1$
structures are fixed independently of the candidate.
\end{proof}

For both logical families, every color is detectable under recoloring, so
$\Det=\Col$, and the number of colors is
\begin{equation}\label{eq:two-parameters}
 N(\Psi_k)=\binom{k}{\lfloor k/2\rfloor}\quad(k\ge2),
 \qquad
 N(\Phi_k)=2^k\quad(k\ge1).
\end{equation}
Each of the structures involved has
\begin{equation}\label{eq:size-threshold}
 n_0(N)=N+1+N\lceil12N^2\ln(2N)\rceil=O(N^3\log N)
\end{equation}
elements: $N$ anchors, the hub, and $N$ cells of $L$ vertices each. Both
values of $N$ in \eqref{eq:two-parameters} are exponential in $k$, and so
is the threshold: $n_0=2^{O(k)}$ for either family. Consequently, nothing
below applies to a target that only has to be correct on structures whose
size is polynomial in $k$, and the construction cannot be adjusted to supply
such structures: property~(a) of Lemma~\ref{lem:palette} already forces
$L\ge N^2+1$ (Remark~\ref{rem:covering}), so structures of the kind used
here have at least $N^3$ vertices however the coloring is chosen.

Both families satisfy \eqref{eq:transport-decoder} once the domain has been
padded with inert anchors, as the next theorem verifies. That theorem and
Corollary~\ref{cor:randomized} below are the precise form of
Theorem~\ref{main:robust}.
\begin{theorem}[Size-specific lower bounds]\label{thm:fixed-size}
Let $\varphi$ be $\Psi_k$ or $\Phi_k$, let $N(\varphi)$ be as in
\eqref{eq:two-parameters}, and let $n\ge n_0(N(\varphi))$ with $n_0$ as in
\eqref{eq:size-threshold}. All the composition-count and term-size lower
bounds of Theorems~\ref{main:omega} and~\ref{main:phi} hold for targets that
are required to be correct only on $n$-element inputs, even when the target
depends arbitrarily on $k$ and $n$.
More generally, let $\Comp$ be an aggregation computation of this Boolean
query in the sense of Section~\ref{sec:responses}, with arbitrary
intermediate scalar values, whose wiring, kernels, scalar constants and
pointwise maps may all have been chosen for $k$ and $n$, and whose primitive
arrays are the inputs of $\varphi$ together with the constant arrays and the
equality array, and nothing else. Then $\wid{\Comp}\ge N(\varphi)$.
\end{theorem}
The last restriction is essential; see the discussion after
Theorem~\ref{thm:transport}.
\begin{proof}
Use the corresponding interpretation of Section~\ref{sec:logical} and add
anchors until the domain has size $n$; this is possible for every $n$ in
the stated range, since the initial domain has $n_0(N)$ elements. The extra
anchors are made inert as follows. For $\Psi_k$, set every entry incident
to an extra anchor to
zero, so that such a witness has both of its edges false at a read-off
pair. For $\Phi_k$, give every entry incident to an extra anchor the code
$\cdb$, so that no type at such an anchor carries a label and such a witness
fails every clause; the core inputs of \eqref{eq:phi-inputs} are kept
unchanged. In both cases the original anchors
and the hub remain excluded as witnesses, for the reasons given in
Section~\ref{sec:logical}, and all the inputs remain type-constant.

It remains to check \eqref{eq:transport-decoder}. For both families a color
is a subset of $[k]$. After recoloring by $f$, a witness in the original
cell $\cell B$ is seen from the hub through $f(B)$ rather than $B$, so for
$\Psi_k$ it satisfies the clauses at $x_A$ exactly when $f(B)\subseteq A$,
hence exactly when $f(B)=A$, and for $\Phi_k$ it succeeds at $x_A$ exactly
when $f(B)=A$. Either way the pair $(x_\lambda,h)$ is satisfied precisely
when $\lambda\in f(\Col)$, which is \eqref{eq:transport-decoder} for every
$\lambda\in\Col$.
Proposition~\ref{prop:fixed-test} now gives the support-number bound. For
\CoR\ we have $\wid{\Comp}=2m$, so $m\ge\lceil N(\varphi)/2\rceil$ since $m$
is an integer, and Equation~\eqref{eq:tree-size} gives the term-size bounds.
\end{proof}
The explicit upper bounds are valid at every size, so the composition count
for $\Phi_k$ remains $\Theta(2^k)$ even when the target may depend on the
size. This is a result about relational circuits, not a lower bound for
unrestricted Boolean circuits with access to individual input entries.

\subsection{A bounded-error consequence}
For fixed $k$ and $n$, a randomized computation is an input-independent
probability distribution over computations in the sense of
Section~\ref{sec:responses}, each of which has its wiring, kernels,
constants and pointwise maps fixed. The distribution may depend on $k$ and
$n$. The error is measured at each input and each output pair. Random
auxiliary relations, random vertex labels, and input-dependent choices of a
circuit are not part of this model: the first two are in general not
type-constant, so Theorem~\ref{thm:transport} does not apply to a
computation that has one of them among its primitives (see the discussion
after that theorem). The distribution may thus range over the computation,
but not over the arrays from which it starts.

\begin{corollary}[Bounded error at a fixed size]\label{cor:randomized}
Let $\varphi$ be either family and $n\ge n_0(N(\varphi))$. A randomized
relational circuit computing $\varphi$ with error at most $\varepsilon<1/2$
on every input and output pair satisfies
\[
 \mathbb E[\gcost{\Comp}]\ \ge\ \frac{1-2\varepsilon}{2}\,N(\varphi).
\]
The same holds for expected composition occurrences in terms. For general
aggregation computations the bound is
$\mathbb E[\wid{\Comp}]\ge(1-2\varepsilon)N(\varphi)$.
\end{corollary}
\begin{proof}
Choose a label $\lambda$ uniformly from $\Col$ and, with equal
probability, test $(x_\lambda,h)$ on $\M$ or on $\M_{f_\lambda}$; the
correct answers are opposite. For a deterministic computation with support
$K$, every $\lambda\notin K$ gives the same output on the two tests, so at
most one of them can be correct, and the success probability is at most
\[
 \frac12+\frac{|K|}{2N(\varphi)}
 \ \le\ \frac12+\frac{\wid{\Comp}}{2N(\varphi)}.
\]
Averaging over the input-independent distribution, a pointwise error of at
most $\varepsilon$ gives an average success probability of at least
$1-\varepsilon$, and comparing this with the bound above yields the claim.
If the expected support number is infinite, the claim is immediate. For
relational circuits, use $\wid{\Comp}=2\gcost{\Comp}$.
\end{proof}
In particular, error $1/3$ still requires $N(\varphi)/6$ compositions in
expectation, that is $2^k/6$ for $\Phi_k$.
